\section{Tokenised single equities}\label{sec:equities}

A tokenised single equity represents a linear, long-only claim on an identified
listed share. We apply the underlying share's SIMM equity delta treatment to that
claim. The settlement assumption is access to the underlying shares and their sale
through an authorised intermediary. The practitioner must verify that the token's
legal rights and custody arrangements permit that sale and payment of the proceeds
to the lending market. A quoted underlying price alone establishes neither access
to the shares nor an executable exit.

The haircut retains the four components of Eq.~\eqref{eq:agg-haircut}. Share sales
use the concentration and clearing-window definitions below. Fund redemptions
retain the definitions of Sections~\ref{sec:theory:capacity}
and~\ref{sec:efficacy}. Leveraged or inverse products, daily-reset funds, options,
and synthetic claims require their own sensitivity and settlement treatment before
admission. The bridge exclusion of Section~\ref{sec:overview-channel} still applies.

\subsection{Underlying share and SIMM bucket}\label{sec:equity-buckets}

The underlying share determines the equity risk factor and bucket. SIMM v2.6
\P\P\,50--56 assign single shares to buckets 1--10 by capitalisation, region, and
sector~\citep{isda_simm_v2_6}. Large capitalisation means at least USD 2 billion,
measured across the issuer's global listings. Buckets 1--4 cover large emerging-market
issuers and buckets 5--8 cover large developed-market issuers. Sector distinguishes
the four buckets within each group. Buckets 9 and 10 cover small emerging-market
and developed-market issuers respectively, across all sectors.

Classification requires evidence of the issuer's most material activity and region.
The practitioner records the underlying identifier and the dated evidence for each
classification input. The developed-market list and the treatment of multinational
issuers follow SIMM \P\P\,53--55~\citep{isda_simm_v2_6}. The exchange on which a
share trades does not by itself establish its region. Missing evidence required for
assignment sends the share to Equity Residual under \P\,56.

Table~\ref{tab:equity-parameters} gives the v2.6 parameters used for single shares.
The risk weight and concentration threshold are taken together from that version.
Bucket 11 retains the fund-unit treatment of Section~\ref{sec:fund-unit}; wrapping
a single share in a token leaves its underlying equity bucket unchanged.

\begin{table}[htbp]
\centering\small
\caption{Single-equity parameters from SIMM v2.6, \P\P\,56, 59 and 77.
Risk weights multiply the sensitivity to a 1\% price move.
Thresholds are USD millions of that sensitivity.}
\label{tab:equity-parameters}
\begin{tabular}{@{}lrrr@{}}
\toprule
Bucket & Risk weight & Correlation & Threshold \\
\midrule
1 & 30 & 0.18 & 3.00 \\
2 & 33 & 0.20 & 3.00 \\
3 & 36 & 0.28 & 3.00 \\
4 & 29 & 0.24 & 3.00 \\
5 & 26 & 0.25 & 12.00 \\
6 & 25 & 0.36 & 12.00 \\
7 & 34 & 0.35 & 12.00 \\
8 & 28 & 0.37 & 12.00 \\
9 & 36 & 0.23 & 0.64 \\
10 & 50 & 0.27 & 0.37 \\
Residual & 50 & 0.00 & 0.37 \\
\bottomrule
\end{tabular}
\end{table}

\subsection{Sensitivity and concentration}\label{sec:equity-concentration}

SIMM equity sensitivity measures the value change for a relative move of one percentage point in the underlying price (\P\,26)~\citep{isda_simm_v2_6}.
For an unleveraged share, the sensitivity per unit of published valuation is
$s_k=0.01$. Thus a risk weight of 50 gives a weighted sensitivity of $0.50$
before concentration. The equivalent exposure-delta convention uses unit delta
with a fractional risk weight of $0.50$. An implementation must convert between
these conventions explicitly when importing Common Risk Interchange Format rows.

The assessed exposure covers the full supply cap and any larger live position.
Let $E_{\mathrm{live}}^{(a)}$ be the current collateral value in USD and define
\begin{equation}\label{eq:equity-exposure}
C^{(a)}=\max\!\left(S_{\max}^{(a)},E_{\mathrm{live}}^{(a)}\right).
\end{equation}
The cap is a proposed input to the assessment. A cap reduction leaves the larger
live exposure in this calculation until that exposure has run off. For one share
denominated in the borrow currency, its USD sensitivity is $0.01C^{(a)}$ and
\begin{equation}\label{eq:equity-cr}
\mathrm{CR}^{(a)}_{\mathrm{EQ}}
=\max\!\left(1,\sqrt{\frac{0.01C^{(a)}}{T_b^{\mathrm{EQ}}}}\right),
\qquad
\mathrm{IM}^{(a)}=\mathrm{RW}_b\,0.01\,\mathrm{CR}^{(a)}_{\mathrm{EQ}}.
\end{equation}
Here $T_b^{\mathrm{EQ}}$ is SIMM's published concentration threshold in USD of
sensitivity. For bucket 10, $T_b^{\mathrm{EQ}}=370{,}000$ USD and concentration
starts above USD 37 million of exposure. The published threshold replaces
$R_{\mathrm{cap}}$ for this sale treatment.

The aggregation equations of Section~\ref{sec:normalisation} apply to the weighted
sensitivities. Distinct equity factors use their own concentration factors, so
$f_{kl}$ can differ from one. Within-bucket correlations follow
Table~\ref{tab:equity-parameters}, and cross-bucket correlations follow SIMM
\P\,60~\citep{isda_simm_v2_6}. Residual remains additive. One share in one bucket
has no pair to diversify, giving Eq.~\eqref{eq:equity-cr}. Any FX sensitivity of a
share denominated outside the borrow currency must also enter the calculation,
with the risk classes assigned to the applicable SIMM product class. The
single-share formula above assumes that FX sensitivity is zero. A currency mismatch
requires the hedge evidence of Section~\ref{sec:tax-market} and an explicit FX
calculation before this simplified formula can support lending limits.

\subsection{Sale horizon and caps}\label{sec:equity-execution}

The sale horizon assumes that the entire assessed exposure liquidates together.
Let $\mathrm{ADV}^{(a)}$ be the underlying share's average daily dollar volume over
a stated number of completed regular trading sessions. The practitioner selects
that observation period and a participation fraction $\pi^{(a)}\in(0,1]$, the
maximum share of daily volume the liquidation may use. These are model choices
registered in Appendix~\ref{sec:free-parameters}. The daily execution capacity and
required number of sessions are
\begin{equation}\label{eq:equity-sessions}
D^{(a)}=\pi^{(a)}\mathrm{ADV}^{(a)},
\qquad
n^{(a)}=\left\lceil\frac{C^{(a)}}{D^{(a)}}\right\rceil.
\end{equation}
The ceiling applies before rounding any displayed volume or cap. The volume
series must identify its source and observation date. Token trading volume and
the issuer's total market capitalisation cannot substitute for the underlying
share's executable volume.

Cash settlement determines the end of the clearing window. Let
$\mathcal{B}^{(a)}(n)$ be the documented worst-case elapsed time in business days
from a liquidation breach to settled cash after $n$ execution sessions. It includes
access to the shares, the wait for an eligible session, and settlement after the
last sale. The relevant exchange and settlement calendars determine that time.

Liquidations outside regular hours wait for the next eligible session. Holidays
and shortened sessions must be reflected in the available execution capacity.
For an ordinary daily calendar with additive access wait $w_{\mathrm{EQ}}$ and
settlement lag $d_{s,\mathrm{EQ}}$, $\mathcal{B}(n)=w_{\mathrm{EQ}}+n+d_{s,\mathrm{EQ}}$.
We set
\begin{equation}\label{eq:equity-lh}
\mathrm{LH}^{(a)}=\max\!\left(10,\mathcal{B}^{(a)}(n^{(a)})\right).
\end{equation}

The minimum of ten business days retains SIMM's native horizon because participation
in observed volume is an execution estimate. It earns no shorter-window margin
credit. The price-risk reserve then follows Eq.~\eqref{eq:m-scaling}, including
its qualified-observation rule. Documented execution fees augment that reserve
under Eq.~\eqref{eq:equity-clearing}. The interest reserve follows
Eq.~\eqref{eq:carry-solve} over the same window.

The practitioner approves a supply cap only after assessing its implied haircut
and lending limits. At that cap the borrow limit must be positive and the haircut
must be below one. The pool-level borrow cap satisfies
\begin{equation}\label{eq:equity-caps}
B_{\max}^{(a)}\le\ell_{\max}^{(a)}S_{\max}^{(a)}.
\end{equation}
The proposed supply cap is held fixed while the calculation runs. Changing the
cap requires a new assessment because both concentration and the sale horizon
depend on it.

Equity feasibility requires evidence that the whole assessed exposure can be sold
and settled within the calculated window. The redemption multiple $M$ and the
one-window redemption-capacity test apply to fund redemptions.
The participation assumption requires validation under stressed volume.

Documented execution fees are an additional cost of settlement. Let
$\epsilon_{\mathrm{exec}}^{(a)}\ge0$ be those fees as a fraction of collateral value,
including any fixed fee divided by the assessed exposure. The equity treatment
extends the price-risk reserve of Eq.~\eqref{eq:m-scaling} to
\begin{equation}\label{eq:equity-clearing}
\begin{aligned}
A^{(a)}_{\mathrm{EQ}}&=\mathrm{IM}^{(a)}\sqrt{\mathrm{LH}^{(a)}/10}
+\epsilon_{\mathrm{exec}}^{(a)},\\
\mathrm{IM}^{(a)}+\delta_{\mathrm{clear}}^{(a)}&=
\begin{cases}
A^{(a)}_{\mathrm{EQ}} & \text{before qualified observations,}\\
\max\!\left(A^{(a)}_{\mathrm{EQ}},\delta_{\mathrm{obs}}^{(a)}\right)
& \text{once }\ge N_{\mathrm{obs}}\text{ qualify.}
\end{cases}
\end{aligned}
\end{equation}
Observed clearing cost includes fees, so the maximum counts them once.
Fee evidence is required before the cap can support new exposure.

Shared execution capacity is counted once across deployments. The practitioner
aggregates exposures to the same underlying share when evaluating concentration
and simultaneous sale demand. Distinct tokens backed by the same share use the
same underlying volume budget. In Eqs.~\eqref{eq:equity-cr} and~\eqref{eq:equity-sessions}, $C^{(a)}$ then denotes the sum of assessed exposures sharing that underlying, with access and settlement assessed for each arrangement.

A fall in volume lengthens the assessed horizon
and requires lower limits for new exposure or a lower cap. Missing or stale volume,
an unavailable sale arrangement, a trading suspension, or an unverified valuation
feed stops new borrowing. Missing execution-capacity evidence also sets the cap
for new collateral inflow to zero. Reductions follow the treatment of existing
positions in Section~\ref{sec:operations}.

\subsection{Valuation and additive layers}\label{sec:equity-layers}

The published valuation follows verified observations of the underlying share.
Between observations Eq.~\eqref{eq:nav-process} applies with $y^{(a)}=0$.
Dividends and corporate actions require an attested adjustment to the token's
share entitlement or cash balance when recognised. A dividend yield is not
deterministic appreciation between price observations. The publication calendar
must specify how the feed treats exchange closures and the next opening price.
The buffer interval $\Delta t_A$ includes those closures and any publication lag.
A frequent feed heartbeat does not shorten an interval with no usable underlying
quotation.

Valuation staleness includes the ordinary gap between an exchange session and
the next usable price. The practitioner sets its baseline and the oracle-integrity
charge under the existing layer definitions. The overdue-publication escalation
still follows Eq.~\eqref{eq:stale}. Wrong-way risk takes a positive charge where
the issuer's business creates a documented connection between the share's loss
and borrower stress, for example a crypto-dependent operating company.

The same fact enters either a named layer or the due-diligence multiplier.
Where the scorecard already charges a specific operational or settlement weakness,
that fact contributes no second charge through $\lambda_{\mathrm{ops}}$ or
$\lambda_{\mathrm{gate}}$. A distinct residual risk requires its own documented
assignment. The private-credit default add-on remains specific to private credit.
We apply the equity risk weight without a separate generic equity jump charge.
That is a model assumption: the square-root rescaling does not establish a bound
on losses from a trading halt, an opening gap, or issuer failure.

\subsection{A small-cap share worked example}\label{sec:equity-example}

A hypothetical developed-market small-cap share falls in bucket 10. Assume
USD 200{,}000 of average daily dollar volume over 21 completed sessions and
10\% participation. Daily execution capacity is USD 20{,}000. The example assumes
one business day of access wait and one business day of settlement after the final
session, with a regular daily calendar. Live exposure is at or below each proposed
cap in Table~\ref{tab:equity-example}.

The example assumes zero incremental execution fees and uses practitioner choices
of $\Lambda=0.025$, $\kappa=0.75$,
$\nu=0.08$, and a ceiling annual borrow rate of $0.20$. These values illustrate
the calculation and are not prescribed SIMM parameters. Every cap shown is below
the USD 37 million concentration threshold, so $\mathrm{CR}=1$ and
$\mathrm{IM}=0.50$. With no qualified clearing observations, the quantities before
the interest reserve and after it are
\begin{equation}\label{eq:equity-example}
H_0=0.50\sqrt{\mathrm{LH}/10}+0.025,\qquad
c=0.20\,\mathrm{LH}/252,\qquad
\theta=\frac{0.75(1-H_0)}{1+0.75c}.
\end{equation}
For an admissible $H_0<1$, the total haircut is $H=H_0+\theta c$ and the borrow
limit is $\max(0,\theta-0.08)$. An $H_0\ge1$ makes the proposed cap ineligible
before the interest calculation, preventing a negative debt or interest reserve.

\begin{table}[htbp]
\centering\small
\caption{Illustrative bucket-10 limits at proposed supply caps. Monetary values
are USD thousands. Haircuts and lending limits are percentages.
The clearing window includes access wait and settlement.}
\label{tab:equity-example}
\begin{tabular}{@{}rrrrrr@{}}
\toprule
Cap & Sessions & $\mathrm{LH}$ & $H$ & LLTV & LTV \\
\midrule
100 & 5 & 10 & 52.78 & 35.41 & 27.41 \\
250 & 13 & 15 & 64.06 & 26.96 & 18.96 \\
500 & 25 & 27 & 84.90 & 11.32 & 3.32 \\
750 & 38 & 40 & Ineligible & 0.00 & 0.00 \\
\bottomrule
\end{tabular}
\end{table}

The sale horizon lowers permitted borrowing well before native SIMM concentration
increases the margin. At the USD 750{,}000 cap, $H_0=1.025$ exceeds the asset's
value and the cap fails the eligibility rule. The example therefore makes supply
capacity an input to lending limits even when the share's SIMM risk weight and
concentration factor remain unchanged.
